\section{总结与延伸}

\subsection{本章小结}

如果把第一讲压缩成一张更实用的清单，答案大概是下面四项：
\begin{enumerate}
    \item 你要能把语言模型视作一个由数据、模型、系统和目标函数组成的完整系统。
    \item 你要知道为什么 tokenization、架构、训练、系统、规模律、数据和对齐不能割裂看。
    \item 你要理解“效率”不是一个口号，而是决定可扩展性和可研究性的硬约束。
    \item 你要开始形成一种研究姿态：先构建、再测量、再放大、最后对齐。
\end{enumerate}

\begin{importantbox}{第一讲的真正交付物}
不是某个模型细节，而是一个工作方式：从抽象入口回到机制层，再把机制放回完整工程栈里重新理解。
\end{importantbox}

\end{document}
